\documentclass[11pt]{article}
\usepackage[T1]{fontenc}
\usepackage[utf8]{inputenc}
\usepackage{lmodern}
\usepackage[a4paper,margin=1in]{geometry}
\usepackage{microtype}
\setlength{\parindent}{0pt}
\setlength{\parskip}{0.5em}

\begin{document}

\begin{center}
{\large\textbf{Executive Summary}}\\[2pt]
{\small Designing, Commercializing and Internationalizing Chinese Spherical Tokamaks --- TUNEM Master's Research Proposal}
\end{center}
\vspace{2pt}

This proposal advances four core claims along one spine --- Capabilities $\rightarrow$ Commercialization $\rightarrow$ International Markets; the first two are technical and form the analytical heart of the thesis.

\textbf{1.~Reconstructing the design interface.} Freidberg, Mangiarotti and Minervini's \emph{Designing a Tokamak Fusion Reactor---How Does Plasma Physics Fit In?} (Phys.\ Plasmas, 2015) shows that a reactor's main parameters are fixed almost entirely by nuclear and engineering constraints, leaving the plasma to satisfy a small set of coupled limits. The thesis derives that ``design interface'' from first principles --- the knobs (field $B$, shape and $\beta$, heating, pulse) and the operational limits (Troyon, Greenwald, kink, bootstrap) that bind them --- rebuilding transparently what the specialist sources compress or assume.

\textbf{2.~Deriving the spherical tokamak as the route that closes Freidberg's constraints.} Freidberg derives that a conventional tokamak does not close: the required plasma current exceeds the kink limit, and the achievable bootstrap fraction ($\approx 0.44$) falls short of the value steady state requires ($\approx 0.84$). One of his own remedies is to raise the magnetic field --- which lowers the required current and eases the kink limit --- and the spherical tokamak reaches that same high-field-favourable regime through \emph{geometry}. Two distinct mechanisms do the work, kept separate: the ST's strong $\tau_E\propto B^{1.0}$ scaling lets confinement come from field rather than current, relieving the kink/current trap, while its low aspect ratio raises the trapped-particle (bootstrap) fraction, relieving the bootstrap shortfall directly; high $\beta$ is then an economic and stability margin, not the cure. The spherical-tokamak reactor literature (Peng; Sykes; Menard) develops low-aspect-ratio designs in the high-bootstrap, high-field regime, but does not engage Freidberg's closure argument; connecting that regime to his specific over-constraint --- showing transparently that the ST relieves the two failure modes he identifies --- is this thesis's own \emph{synthesis}, accessible exposition rather than new physics.

\textbf{3.~Commercialization --- the core product problem.} The heart of commercializing a reactor is making a useful product people will pay for, and \emph{scientific} breakeven is not economic viability. The thesis derives that gap --- separating what the spherical tokamak has \emph{demonstrated} (high $\beta$, bootstrap fraction, compactness) from the hard walls that \emph{remain} (sustained high-field confinement, tritium breeding in cramped geometry, neutron-tolerant materials, continuous operation) --- and frames it as the directed engineering agenda the entrepreneur must close before a commercial product exists. Capital, regulation, supply chain and the other hurdles of building a fusion company are taken up only \emph{after} that core product problem, at literature-review depth on public sources.

\textbf{4.~Internationalization --- taking the technology abroad.} The thesis closes by asking what deal structures (turnkey loans, equity stakes, joint ventures) could carry a Chinese fusion technology to partner countries, with the Pakistan/Argentina contrast as the evidence base. It weighs one open question both ways: whether fusion's lack of fissile fuel and weapons-usable plutonium outweighs the residual tritium-breeding and neutron-activation concerns enough to ease the sovereignty anxiety that stalled fission exports such as Atucha~III. This integrated synthesis --- absent from a literature that is either specialist or business-only --- is what makes TUNEM's engineering-and-management mission concrete as an actionable blueprint for the engineer-entrepreneur, and it is the foundation of the author's intended post-doctoral ventures.

\textbf{Primary sources.} The analytical core rests on Freidberg, Mangiarotti \& Minervini (2015); Freidberg, \emph{Plasma Physics and Fusion Energy}; Wesson, \emph{Tokamaks}; the spherical-tokamak reactor literature (Peng; Sykes; Menard); and Prof.\ Tan Yi's SUNIST-2 programme, whose stated aims --- confinement at higher \emph{magnetic field}, magnetic-reconnection heating, plasma-shaping and control, and repetitive pulsed operation --- map directly onto the design knobs, the field knob above all.

\end{document}
